\section{Method}
\label{sec:method}

The dual probe operationalises stale coherence by measuring, at the
same simulation instant, the fraction of injected anchors that the
model echoes in its verbal output and the fraction whose
corresponding action keys are weighted in its action output. To
make the probe meaningful, the substrate has to expose a verbal
channel and an action channel that share a single underlying model
state. This section describes the substrate, the anchor injection
protocol that prepares the probe surface, and the probe definitions
themselves.

\subsection{Multi-Cadence Structured-Output Substrate}
\label{sec:method:substrate}

We instantiate the dual probe on a deterministic two-dimensional
colony simulation. A small group of worker units gather food,
wood, and stone from tiles on a tile-world map, deposit them in
storage, construct simple structures (shelter, walls, watchpoints),
and respond to periodic threats (weather shifts, raids, faction
tension). There is no terminal win condition; the simulation
persists for the configured horizon as long as the colony's
population threshold survives, which it always does in the no-LLM
control. Our purpose is not to score whether the LLM keeps the
colony alive---under the model families we test it does---but to
expose a setting in which a single LLM emits parallel verbal and
action directives over a sustained horizon against a fixed
deterministic tool surface (A* path planning, Boids movement,
seeded RNG, a build planner).

A single LLM drives the colony across four channels.
An \emph{environment-director} emits weather, event, and
faction-tension shifts every 5--15 simulation seconds; an
\emph{npc-policy} emits per-group target priorities and intent
weights on the same fast cadence; a \emph{strategic-plan} emits
prose summaries and high-level goals every 90 simulation seconds;
and a \emph{colony-agent} emits per-step build plans through a
Perceive--Plan--Ground--Execute--Evaluate--Reflect cycle on a
slower sim-time gate. Each channel emits JSON that is validated
against a typed response schema and clamped to feasible numeric
ranges before it acts on the world. The probe surface is the
cross-product of two channels: the strategic-plan's prose carries
the verbal signal, and the npc-policy's target priorities and
intent weights carry the action signal. Table~\ref{tab:channels}
summarises the assignment.

The multi-cadence emission shape is the structural property that
makes verbal--action dissociation detectable. A single-cadence
harness conflates the verbal summary and the next action into one
inference call, so they cannot drift. By emitting strategic prose
on a slow cadence and policy-weight vectors on a fast cadence over
the same shared memory, the substrate exposes a verbal channel and
an action channel that can move independently.

\begin{table}[t]
  \centering
  \footnotesize
  \begin{tabular}{l l l}
    \toprule
    Channel & Cadence & Probe surface \\
    \midrule
    environment-director & 5--15\,s & --- \\
    npc-policy           & 5--15\,s & action \\
    strategic-plan       & 90\,s    & verbal \\
    colony-agent         & sim-gated & --- \\
    \bottomrule
  \end{tabular}
  \caption{The four channels emitted by the substrate, their
  cadences, and which probe surface each contributes to. The
  environment-director and colony-agent channels are run for
  realism but are not part of the verbal--action recall probes.}
  \label{tab:channels}
\end{table}

\paragraph{Why a custom substrate.}
Existing agent benchmarks do not expose the emission shape the
dual probe requires. Crafter \citep{Hafner2022Crafter} and Voyager
\citep{Wang2023Voyager} place a single LLM-as-policy in a tile or
voxel world but emit one action per step, with no parallel
structured directives across cadences; verbal and action cannot
dissociate because they share an inference call. AgentBench
\citep{Liu2024AgentBench} and tool-use benchmarks evaluate
task-discrete episodes that reset the context each run,
eliminating the sustained joint emission the dual probe requires.
BALROG \citep{Paglieri2025BALROG} stresses agentic reasoning on
NetHack-style games but uses a single textual action channel.
Cicero \citep{FAIR2022Cicero} does emit parallel dialogue and
action streams---and is the closest existing artifact to what we
need---but is closed-source and not configurable for controlled
anchor injection. Custom-building the substrate gives us the four
properties the dual probe requires (a verbal channel, an action
channel, multiple cadences, a deterministic tool surface), and an
anchor injection point, without entangling them with the
proprietary specifics of a production deployment.

\paragraph{Action chain.}
The substrate closes the loop between LLM directive and world
state. On each new npc-policy envelope, every worker's FSM state
is re-sampled from \textsf{intent\_weights}; its target tile is
the highest-weighted \textsf{target\_priorities} key compatible
with that state (harvest states resolve to farm/lumber/quarry
anchors; deliver states resolve to warehouse). Harvest-state
workers that reach their tile add food, wood, or stone to the
colony stockpile each tick at a calibrated yield, while an
independent resource system depletes food proportional to alive
worker count regardless of LLM input. A directive that
under-weights harvest intents or biases targets toward
deposit-only anchors visibly bleeds the food reserve, and the
updated counts enter the next prompt's world summary.

\paragraph{Determinism and balance.}
The substrate is deterministic in the no-LLM control: identical
seeds reproduce identical world states bit-for-bit across runs and
operating systems. Determinism is enforced through a single seeded
PCG64 generator with sub-streams derived by keyed BLAKE2b, a
versioned numpy tile grid, and deterministic implementations of
A* path planning and Boids movement. When the LLM is active, only
its sampling introduces stochasticity; everything downstream of an
emitted directive is deterministic. The $|\Delta| = 0.3$ threshold
is calibrated against this control: with the LLM replaced by a
stub adapter that emits constant policy weights, the dual probe
returns $\Delta \approx 0$ with a $|\Delta|$ noise floor below
$0.05$ across all seeds, defining what counts as a substantive
signal rather than substrate jitter. Game rules---resource gather
rates, structure costs, threat frequencies---are constants from
the simulation's config, deterministic given the seed, and
identical for every LLM family the pilot evaluates; the substrate
is the same lab for every model and our claims are differential.

\subsection{Anchor Injection Protocol}
\label{sec:method:anchors}

At simulation time $t = 0$ we inject five anchor entries into the
LLM's persistent world summary. Each anchor is a paired tuple of a
verbal token and an implicit goal:

\begin{equation*}
  a_i = (\text{verbal\_token}_i,\ \text{implicit\_goal}_i).
\end{equation*}

A representative anchor is
$(\textsf{warehouse at (12,8)},\ \{\textsf{action}: \textsf{deliver},
\textsf{target}: \textsf{warehouse}\})$. The two halves of the
anchor are co-located in the world-summary memory but are not
presented to the model as an instruction. The model must
autonomously integrate the pair: the verbal half may surface in its
strategic-plan output (verbal recall) and the action half may
surface in its npc-policy weights (action-grounded recall). This
makes verbal and action recall independently checkable on the same
underlying anchor set.

The five-anchor count is calibrated against pilot runs we conducted
during substrate development. Below five anchors the verbal probe
saturated near $1.0$ for all three families, eliminating its
discriminative power; above five anchors smaller open-weight models
collapsed their structured-output validity rate. Five anchors
preserved both the discrimination ceiling of the verbal probe and
the JSON validity floor of the weakest tested family. The number
of channels (four), the cadence stratification (5--15\,s for fast
events, 90\,s for strategic prose, sim-gated for the colony loop),
and the JSON schema validation are similarly load-bearing for the
probe: collapsing the cadence stratification to a single rate
removes the dissociation surface (we ran single-cadence variants
during design and obtained $\Delta \approx 0$ across families), and
removing schema validation breaks the action-key vocabulary the
dual probe scores against. The remaining substrate
details---tile-world dimensions, worker counts, structure
costs---are mundane: they exist so the LLM has something concrete
to operate over, but they are not load-bearing for the probe and
can be re-instantiated in any other domain that exposes the same
four properties.

\subsection{The Dual Probe}
\label{sec:method:probe}

The dual probe samples two quantities at a fixed sim-time $t$.
\emph{Anchored fact recall} (AFR) is the fraction of the five
injected anchors whose verbal token appears in the LLM's verbal
output blob aggregated across the strategic-plan, environment
director, per-group policy summary, and colony-agent rationale
fields at time $t$:

\begin{equation*}
  \mathrm{AFR}(t) = \frac{1}{|A|} \sum_{a_i \in A} \mathbf{1}\bigl[
    \text{verbal\_token}_i \in \mathcal{V}(t) \bigr].
\end{equation*}

\emph{Action-grounded recall} (AGR) is the fraction of the five
anchors whose implicit-goal action key appears in any group's
\textsf{target\_priorities} or \textsf{intent\_weights} with weight
at least $\tau = 0.3$ on the clamped $[0, 3]$ scale used by the
guardrails:

\begin{equation*}
  \mathrm{AGR}(t) = \frac{1}{|A|} \sum_{a_i \in A} \mathbf{1}\bigl[
    \exists g:\ w_g(\text{action\_key}_i; t) \geq \tau \bigr].
\end{equation*}

The dissociation gap is $\Delta(t) = \mathrm{AFR}(t) -
\mathrm{AGR}(t)$. The signed gap distinguishes two directions of
dissociation. $\Delta > 0$ corresponds to stale coherence, in which
the model continues to mention the anchor while the action channel
no longer weights it. $\Delta < 0$ corresponds to the inverse
pattern, in which the model has dropped the anchor from its verbal
output while the action channel still weights it. We report the
signed gap and the magnitude $|\Delta|$ against a pre-registered
threshold of $0.3$.

Stale coherence as a deployment-relevant phenomenon requires more
than a non-zero $\Delta$: the action channel must also continue to
act on the world rather than collapsing into a no-op. We discount
the gap by a behavioural-drift measure, computed as the Jaccard
distance between the set of weighted action keys at $t = 0$ and at
$t$. When drift is zero the policy has remained at its initial
configuration regardless of input, and a non-zero $\Delta$ is not
informative; we exclude such cells from the threshold-crossing
count and report them transparently in the per-cell table. The
drift discount disentangles dissociation (the model is acting but
on a different anchor distribution than its prose describes) from
goal displacement (the model has redirected the policy elsewhere)
and from no-op output (the model is not acting at all).
